\subsection{Source Attribution at the Close of a 30-s Window}
\label{sec:window-source-localization}
Can the monitor still assign a source when the dynamic bank rejects the
startup noise model? We tested a window-level fallback in a separate
IEEE 39-bus simulator that superimposes prescribed disturbances on normal
phasor trajectories. Each record begins with 60\,s of quiet measurements.
Both procedures see the same eight PMUs and report once, at the end of
the next 30\,s. This deadline permits use of the entire transient,
including post-event measurements. The timely-attribution experiment in
the main article requires a diagnosis before restoration.

The comparator starts with the competition classifier and checks for
missing reports and phase imbalance. When the dynamic bank rejects the
angular noise scale, the comparator scores ten-second response templates
for candidate assets. It standardizes voltage and frequency channels
against the quiet prefix, omits missing observations and fits a positive
response amplitude at each candidate onset. Its calibration threshold
exceeds the largest score on development windows containing only normal
operation or isolated instrument disturbances. System-wide frequency
helps identify the event family, but a frequency change shared by all
PMUs contributes little information about which bus switched.

The revised procedure estimates the source within the predicted family
using the voltage and current magnitudes and angles across PMUs. For a
persistent generation or load change, the observed displacement $z$ is
the mean over the last ten seconds relative to the quiet prefix. For
each candidate $c$, let $s_c$ be its mean displacement in development
records and $W$ the inverse pooled within-family covariance after
standardization and shrinkage. All quantities are restricted to
observed coordinates before covariance inversion. The source estimate is
\[
   \widehat c=\arg\min_c\ \min_{0.4\leq a\leq2}
         (z-a s_c)^\top W(z-a s_c).
\]
The amplitude interval is relative to the development signature. This
fit allows an intervention to differ in size while preserving the
spatial pattern that distinguishes two candidates.

The ninth template pairs a weak physical pulse with a current gain
change at one PMU and a simultaneous change of angle encoding at
another. Adding a full turn leaves the electrical phasor unchanged.
The monitor uses the encoding change as an instrument flag and uses
wrapped angles for the source estimate. It requires a three-phase
median current-magnitude excursion above 12\% and a voltage-angle
excursion more than $90^\circ$ beyond the prefix range on at least two
phases of the other PMU. Each flag must last at least 12 reports, with
at least nine reports of overlap. To localize the physical pulse, it
fits the quiet oscillations outside the affected interval using the two
dominant low-frequency components of startup voltage magnitude. It
subtracts this baseline over the pulse and omits the contaminated
current magnitude. Neither affected PMU is supplied as the physical
source. This rule recognizes a known composite template.

We fixed the source signatures and decision rules before testing.
Development comprised 216 balanced records and 150 additional composite
records. The primary test included 24 new records per class. Each
physical family had three candidate locations. We shuffled their order
and varied noise, onset, severity and duration. Instrument faults
rotated across PMUs. The 24 normal windows and 48 further normal records
provide 0.6 monitored hours for the false-alarm comparison. Source
accuracy uses all 192 nonnormal windows, including missed events.
In this simulator the response attenuates with distance in bus
numbering. The resulting source discrimination is specific to its
prescribed spatial signatures.

\begin{table}[t]\centering\small
\caption{Paired nine-class window decisions on independently generated
records. The last row counts normal windows containing an alarm.}
\label{tab:window-source-comparison}
\begin{tabular}{lrr}\toprule
Measure & Response bank & Revised procedure\\\midrule
Macro-F1 & 0.845 & 0.991\\
Correct family & 190/216 & 214/216\\
Exact source & 166/192 & 188/192\\
Alarms on normal windows & 0/72 & 0/72\\\bottomrule
\end{tabular}
\end{table}

All 264 records failed the original dynamic-bank angle-noise check.
The revised fallback recognized all 24 composite windows and located
the physical source in 22. Both procedures identified the source in
every line, generation and load window in this test. The revised
procedure's two family errors were isolated instrument corruptions
classified as normal. Stratified bootstrap resampling of whole test
records within class (4,000 draws) gave 95\% intervals of
$[0.977,1.000]$ for macro-F1, $[0.977,1.000]$ for correct family and
$[0.958,0.995]$ for exact source. Zero alarms in 72 independent normal
windows gives a one-sided 95\% upper bound of 4.1\% per window under
the tested conditions.

We also evaluated the composite flag after mapping all voltage angles
to $[-180^\circ,180^\circ)$. This transformation preserves every
electrical phasor. The flag fired on 24 of 24 composite records with
the original encoding and on none after conversion. This diagnostic
isolates the flag; it does not rescore the full classifier. The
window-level macro-F1 therefore depends on an instrument convention
in these records.

\paragraph{Transfer to held-out physical sites.}
A second test reserved three new physical assets per family and kept the
same eight PMUs. Faults, outages, generation changes and load changes
were applied at sites outside their respective development catalogs.
The physical components of the two mixed classes also moved. The
instrument-only classes reused the installed PMUs. Each of the nine
classes contributed 12 new records, with 48 additional normal records
reserved for counting false alarms. We compared the submitted competition
classifier, the spatial procedure above, and a version that replaces the
whole-turn flag and maps angles to their signed principal range.

For each phase at each PMU, we fit a local equivalent admittance on
the valid quiet samples $Q$ and measured the subsequent current mismatch,
\begin{equation}
 \widehat Y=\frac{\sum_{t\in Q}\overline{V_t} I_t}
                         {\sum_{t\in Q}|V_t|^2},\qquad
 e_t=\frac{|I_t-\widehat Y V_t|}{\operatorname{median}_{u\in Q}|I_u|}.
\end{equation}
We take the median mismatch across phases at each PMU and subtract the
contemporaneous median across available PMUs. The detector requires a
local excess lasting 12 to 150 reports above both 0.12 and three times
its quiet 99.5th percentile. This residual measures departure from the
local quiet voltage--current relation. A current gain change alone can
produce such a departure, so the residual by itself cannot identify a
simultaneous physical disturbance. Both hybrids use the same source
templates. We sealed the three methods' predictions before scoring.

\begin{table}[!ht]
\centering\small
\caption{Paired window decisions at held-out physical sites. The source
denominator includes all nonnormal windows, including missed events.}
\label{tab:window-held-locations}
\begin{tabular}{lrrr}\toprule
Measure & Submitted & Spatial & Residual\\\midrule
Macro-F1 & 0.193 & 0.991 & 0.991\\
Correct family & 32/108 & 107/108 & 107/108\\
Exact source & 0/96 & 23/96 & 23/96\\
Normal alarms & 0/60 & 0/60 & 0/60\\\bottomrule
\end{tabular}
\end{table}

All 156 records failed the dynamic bank's startup angle-noise check.
Both window hybrids identified the family in all 72 records with a
physical intervention and located none of those sources. Their 23
source hits came from the 24 instrument-only events. The residual
hybrid gave the same family and source on all 156 windows when one
turn was added to every reported angle and the window pipeline was
rerun with the same startup decision. The 60 normal windows represent
30 minutes of observation. These results concern synthetic site
transfer within this simulator; the competition's hidden test records
were unavailable.

\paragraph{Electrical source signatures at previously unused assets.}
We froze the residual classifier and changed only its source ranking.
For a network admittance $Y$ augmented by generator Norton equivalents
with assumed transient reactance 0.2 pu, a current change at bus $k$
produces a voltage signature proportional to $Y^{-1}e_k$.
Opening branch $b$, with admittance stamp $Y_b$, instead gives
$(Y-Y_b)\Delta V=Y_b V_0$ at the archived operating phasor $V_0$.
We restrict each signature to the eight PMUs and divide each voltage
change by its operating phasor. For each candidate, weighted least
squares fits an unknown complex strength and a common complex offset
to the observed fractional voltage changes. The quiet prefix supplies
the inverse-variance weights. The source search includes all 39 buses
for faults and weak composite disturbances, all 46 branches for
outages, the ten generator terminals for generation changes, and the
18 buses with positive active load for load changes. For the mixed
generation and missing-report class we search the generator terminals.

A third cohort used 12 records per class and 48 additional normal
windows. Its physical targets were absent from both earlier site
lists. All three methods saw the same measurements, and their
predictions were sealed before scoring. The simulator imposes spatial
attenuation by numerical bus-label proximity, whereas the electrical
localizer predicts attenuation and phase from network admittance.
This comparison measures transfer across that model discrepancy.

\begin{table}[!ht]
\centering\small
\caption{Nine-class site transfer with a network-derived source dictionary.
All normal-alarm counts cover the same 30 minutes.}
\label{tab:window-electrical-sites}
\begin{tabular}{lrrr}\toprule
Measure & Submitted & Residual & Electrical\\\midrule
Macro-F1 & 0.193 & 0.981 & 0.981\\
Correct family & 32/108 & 106/108 & 106/108\\
Exact source, all events & 0/96 & 22/96 & 24/96\\
Exact physical source & 0/72 & 0/72 & 2/72\\
Normal alarms & 0/60 & 0/60 & 0/60\\\bottomrule
\end{tabular}
\end{table}

All three true assets per physical family were eligible candidates.
Both hybrids classified all 72 physical interventions correctly.
Electrical ranking recovered one line outage and one load change.
The remaining 22 source hits were instrument-only events. All 156
records failed dynamic-bank startup. The static voltage dictionary
also has intrinsic ambiguities: after fitting the common offset,
13 pairs among the 39 bus signatures are proportional at the observed
PMUs (absolute normalized inner product within $10^{-10}$ of unity).
With unknown complex strength, each such pair gives the same
noise-free fit. This ambiguity persists even with an exact network
model; the test does not separate its contribution to the observed
errors from simulator mismatch.
